% chapters/25_video_models.tex
% Part of: AI / LLM / Agents / Hardware Notes

\section{Video \& Predictive Models}

\subsection{Sora \& video diffusion}

Sora (OpenAI, 2024) represents a major milestone in generative video: a text-to-video foundation model capable of producing coherent videos up to 1 minute in length at 1080p resolution. Unlike earlier text-to-video systems that required single-frame generation followed by upsampling, Sora generates video natively and directly.

\subsubsection{Architecture: Video Diffusion Transformer}

The core innovation is the \textbf{Video Diffusion Transformer (DiT)}, which departs from the U-Net-based architecture standard in image diffusion. Rather than applying convolutional layers hierarchically, Sora's transformer operates on \emph{spacetime patches}—small cuboid regions of video (e.g., a $16 \times 16 \times 2$ block spanning spatial and temporal dimensions). Each patch is tokenised into a flat embedding sequence, allowing the transformer to treat video generation as a sequence-to-sequence task similar to language modelling.

This patch-based representation is elegantly general: instead of hardcoding separate models for different resolutions or aspect ratios, Sora simply conditions on the target resolution and generates patches of the appropriate size and duration. A $480p$ portrait video, a $1080p$ landscape shot, and a $720p$ square clip can all be generated by the same model by adjusting patch scaling.

\begin{intuitionbox}
Why patches instead of pixels? Pixel-level diffusion on video is prohibitively expensive: a 1-minute 1080p video at 30fps is $\approx 1.8M$ frames of $1920 \times 1080$ pixels each. Patches reduce this to $\sim 10k$ tokens, making transformer-scale generation tractable. The trade-off is that patch-level generation may lose fine-grained spatial detail, but token efficiency wins.
\end{intuitionbox}

\subsubsection{Training and Data}

Sora was trained on large-scale video datasets paired with text descriptions. A critical practical detail: raw video captions (e.g., YouTube metadata) are often sparse or vague. OpenAI likely employed a \textbf{recaptioning pipeline}: passing each video through a vision-language model (similar to CLIP or multimodal GPT-4V) to generate rich, detailed text descriptions that better explain the visual content. This densification of captions significantly improves text-to-video coherence.

The model learns in the diffusion framework: corrupted videos (with noise added to patches) are reconstructed conditioned on text embeddings. As training progresses, the model learns to generate spatially coherent frames and temporally consistent motion.

\subsubsection{World Simulator Hypothesis}

The Sora paper makes a bold claim: by learning to predict the next patch in a sequence, the model implicitly learns a \emph{world model}—an internal representation of physics, causality, and object dynamics. Evidence cited includes:
\begin{itemize}
  \item Objects maintain consistent identity and trajectories across frames.
  \item Shadows, reflections, and lighting respond realistically to object motion.
  \item Fluid and deformable objects (smoke, water, cloth) obey rough physical laws.
\end{itemize}

This hypothesis is influential but contested. Sceptics argue that the model may simply memorise large-scale patterns from training data rather than developing a true causal model of physics. Distinguishing learned dynamics from pattern matching remains an open question in interpretability.

\subsubsection{Known Limitations}

As of early 2024, Sora exhibits several failure modes:
\begin{warnbox}
\textbf{Physics violations}: Objects can appear or disappear without physical justification. A person may phase through a wall.

\textbf{Long-range consistency}: Over 30+ seconds, character appearance, lighting, and background slowly drift or degrade.

\textbf{Control granularity}: Fine-grained directives (e.g., "move the camera exactly 30 degrees clockwise") are difficult; the model interprets text-level intent rather than precise camera parameters.

\textbf{Temporal flicker}: Despite transformer-based processing, occasional frame-to-frame jitter occurs, especially at patch boundaries.
\end{warnbox}

These limitations suggest that the current generation of diffusion-based video models, while impressive, remain far from true physical simulation or precise interactive control.

\subsection{Genie \& interactive world models}

Genie (Google DeepMind, 2024) tackles a different problem: learning a \emph{controllable} world model from unlabelled video. Rather than generating videos from text prompts, Genie learns to generate video frames conditioned on learned \emph{actions}, enabling interactive control without a pre-defined control scheme.

\subsubsection{Unsupervised Action Learning}

The key innovation is that Genie learns its action space \textbf{without any action labels}. Given a sequence of raw video frames (e.g., from gameplay footage, dashcam footage, or other dynamic scenes), Genie infers: what latent factors, if manipulated, would steer the world's evolution?

Formally, Genie models $P(\text{frame}_{t+1} | \text{frames}_{1:t}, \mathbf{a}_t)$, where $\mathbf{a}_t$ is a learned latent action token. By training this mapping end-to-end on video sequences, the model discovers an action representation where nearby latent tokens correspond to coherent world perturbations (e.g., "move forward," "jump," "rotate camera"). This is enabled by the structure of the data itself: video sequences naturally contain the signal of causality (action causes state change).

\begin{intuitionbox}
Why unsupervised? Collecting action labels (a button press or joystick input for every frame) is expensive and domain-specific. Gameplay footage, for instance, may come without controller logs. Genie's insight is that the video data itself encodes action information through the sequential structure: if frames are statistically different in a consistent direction, an action likely caused that difference.
\end{intuitionbox}

\subsubsection{Architecture and Components}

Genie's pipeline comprises three main components:

\begin{enumerate}
  \item \textbf{Spatial Tokeniser}: A VQ-VAE (Vector Quantised VAE) encodes individual video frames into discrete latent codes. This compresses spatial information and prepares frames for the dynamics model.

  \item \textbf{Latent Action Model}: A learned discrete action tokeniser discovers a set of "primitive actions" such that sequencing them produces coherent trajectories. This acts like learning a proto-language of control.

  \item \textbf{Dynamics Model}: A diffusion-based or autoregressive model predicts the next frame's latent code given the previous frame and a latent action token. This is the core world model: $P(\mathbf{z}_{t+1} | \mathbf{z}_{1:t}, \mathbf{a}_t)$ where $\mathbf{z}$ are frame latents.
\end{enumerate}

Once trained, a user can "play" the world by selecting action tokens; the model generates a stream of coherent frames responding to those actions. Notably, no traditional controller input is required—the model infers control directly from user interaction (e.g., keyboard or learned policy).

\subsubsection{Connection to World Models}

Genie embodies a modern realisation of the \emph{world models} programme from reinforcement learning: the hypothesis that agents learn an internal forward model of environment dynamics, which they use for planning. Unlike classical world models trained on RL environments with action labels, Genie learns from passive observation—the same regime humans often learn from (watching video, reading books). This makes it potentially more data-efficient and more aligned with human-like learning.

\begin{paperbox}[title={Relation to V-JEPA and Self-Supervised Video}]
Genie and V-JEPA (Chapter 24) are complementary: V-JEPA learns a high-level latent representation of video through masked prediction; Genie learns \emph{actionable} representations where latent variables map to world state changes. Both leverage the structure of video to extract generalizable representations without labels.
\end{paperbox}

\subsection{Temporal consistency challenges}

Video generation introduces challenges absent from image generation: frames must cohere across time, objects must persist and move realistically, and long-range dependencies must be maintained. This subsection surveys the main challenges and techniques.

\subsubsection{Flickering and Inter-Frame Jitter}

Diffusion models, even transformer-based ones, can process frames with too much independence, leading to:
\begin{itemize}
  \item \textbf{Flicker}: Pixel values oscillate between frames without clear motion, resembling poorly denoised noise.
  \item \textbf{Jitter}: Objects appear to vibrate or twitch due to slight spatial shifts in consecutive frames.
\end{itemize}

Root cause: the noise prediction objective optimises each timestep independently; without explicit temporal coupling, the model may find local minima that improve a single frame at the cost of inter-frame coherence.

\subsubsection{Temporal Modelling Approaches}

Several techniques enforce temporal consistency:

\begin{enumerate}
  \item \textbf{3D Convolutions}: Inflate standard 2D convolutions into 3D by adding a temporal dimension. A $3 \times 3$ kernel becomes $3 \times 3 \times 3$, allowing the model to reason about local motion patterns. U-Net diffusion backbones can be retrofitted with 3D convs in the deeper layers.

  \item \textbf{Temporal Attention}: Add self-attention across the time dimension in addition to spatial attention. A transformer can attend to "where was this object in the previous frame?" explicitly, guiding consistency. This is computationally expensive (quadratic in sequence length) but highly effective.

  \item \textbf{Causal Temporal Transformers}: Structure the transformer so that frame $t$ can attend to frames $0 \ldots t-1$ but not future frames. This respects causality and is more efficient than full bidirectional attention, though models that attend bidirectionally before generation (and then use causal generation) are also common.

  \item \textbf{Optical Flow Supervision}: Use optical flow (pre-computed or predicted) as a conditioning signal or auxiliary loss. If the model predicts inconsistent flow between frames, a loss penalty encourages it to enforce motion consistency. Some models compute optical flow between diffusion steps and use it to warp frames for better coherence.
\end{enumerate}

\subsubsection{Motion Blur, Ghosting, and Trajectory Tracking}

Generating realistic motion is non-trivial. The model must:
\begin{itemize}
  \item Predict where objects will be several frames hence.
  \item Render motion blur or sharp motion as appropriate to the scene context.
  \item Avoid \textbf{ghosting}: trailing artifacts where an object's old and new position both appear, creating a "smeared" look.
\end{itemize}

Successful motion generation requires the model to implicitly track object trajectories. Models trained with temporal attention or 3D convolutions tend to perform better because they can reason about object identity across frames.

\subsubsection{Long-Range Consistency}

A harder problem than frame-to-frame coherence: maintaining consistency over hundreds of frames. Issues include:
\begin{warnbox}
\textbf{Character drift}: A person's appearance (clothing, hairstyle) may subtly shift over 30+ seconds.

\textbf{Lighting creep}: The ambient light colour and intensity may slowly change without motivation.

\textbf{Background accumulation}: Repeated small errors in background rendering accumulate, eventually making the scene look wholly different.

\textbf{Layout collapse}: The spatial relationships between objects become confused, leading to impossible geometry.
\end{warnbox}

Current models degrade noticeably after $\sim 1$ minute, suggesting that purely feed-forward generation (frame-by-frame without global memory) is insufficient. Research directions include hierarchical generation (high-level scene layout first, then detail infilling) and memory-augmented architectures.

\subsubsection{Optical Flow as a Supervision Signal}

Some recent models use \textbf{optical flow} (dense motion estimation) as auxiliary supervision:
\begin{itemize}
  \item \textbf{Flow prediction loss}: In addition to the diffusion loss, encourage the model's latent representations to predict accurate optical flow.
  \item \textbf{Flow-based warping}: Use predicted flow to warp frames from frame $t$ to frame $t+1$, then blend the warp with the generated frame. This enforces physical continuity.
  \item \textbf{Flow conditioning}: Condition the diffusion model on precomputed or user-specified optical flow, constraining motion to be plausible.
\end{itemize}

This approach leverages classical computer vision (optical flow estimation) to ground modern generative models, improving consistency at a modest computational cost.

\subsection{Next-frame prediction as pretraining}

Predicting the next frame in a video sequence is a powerful self-supervised objective. It forces the model to learn dynamics, object persistence, and causal structure—all without labels. This subsection covers video prediction as a pretraining paradigm.

\subsubsection{Video Prediction as Self-Supervised Learning}

The core insight: a model that can predict $\text{frame}_{t+1}$ from $\text{frames}_{1:t}$ must understand:
\begin{itemize}
  \item \textbf{Motion}: How do objects move in this domain?
  \item \textbf{Object permanence}: Does an occluded object continue to exist?
  \item \textbf{Physics}: What are the constraints on how the world can change?
  \item \textbf{Semantics}: What is the scene about, and what are the key actors?
\end{itemize}

Unlike supervised pretraining (contrastive learning, masked language modelling), video prediction is a \emph{generative} task: the model must synthesise the raw sensory signal of the next moment. This is a rich learning signal that captures the structure of the world.

\begin{paperbox}[title={Historical Roots}]
Video prediction dates back to the early 2010s (Zeiler et al., Srivastava et al.), where convolutional models learned to predict pixel space. The field has evolved through adversarial losses (GANs), flow-based methods, VAEs, and now diffusion. The core idea remains: learning from temporal structure.
\end{paperbox}

\subsubsection{VideoMAE: Masked Autoencoding for Video}

VideoMAE (He et al., Meta, 2022) applies the \textbf{masked autoencoder} paradigm to video:
\begin{itemize}
  \item \textbf{Masking}: Randomly mask $90\%+$ of spacetime patches in a video. This is much more aggressive than image masking ($\sim 75\%)$ and tests the model's ability to reason about long-range structure.
  \item \textbf{Reconstruction}: An encoder embeds the visible patches; a decoder reconstructs the masked patches. The loss is the reconstruction error on masked regions.
  \item \textbf{Learned representations}: The trained encoder produces video embeddings useful for downstream tasks (action recognition, video classification) without finetuning.
\end{itemize}

VideoMAE demonstrates that masked prediction scales well: larger models, more data, and higher masking ratios all improve downstream performance, following predictable scaling laws. This makes video prediction an attractive foundation for video understanding.

\subsubsection{V-JEPA: Prediction in Representation Space}

Rather than predicting pixels, \textbf{V-JEPA} (Video Joint-Embedding Predictive Architecture, Chapter 24) predicts in a learned latent space:
\begin{itemize}
  \item \textbf{Representation learning}: An encoder $f$ maps frames to latents: $\mathbf{z}_t = f(\text{frame}_t)$.
  \item \textbf{Latent prediction}: A predictor $g$ is trained to minimise the distance between predicted latents $g(\mathbf{z}_{1:t})$ and target latents $f(\text{frame}_{t+1})$ (using a momentum-updated target encoder, similar to BYOL).
  \item \textbf{Advantage over pixel prediction}: Latent space is much simpler than pixel space; the model avoids modelling irrelevant textural variation and learns more abstract motion representations.
\end{itemize}

V-JEPA achieves strong video understanding performance and is computationally cheaper than pixel-space prediction. The representation learned this way captures semantically meaningful motion rather than superficial pixel transitions.

\subsubsection{Scaling Laws in Video Pretraining}

Like language and image models, video pretraining exhibits \textbf{scaling laws}:
\begin{itemize}
  \item \textbf{Data scaling}: More video data yields better performance, with diminishing returns following a power law (similar to language models).
  \item \textbf{Compute scaling}: Larger models achieve lower prediction loss, consistent with Chinchilla scaling.
  \item \textbf{Task transfer}: Models trained on video prediction with internet-scale data transfer to many downstream tasks (classification, detection, tracking) better than supervised pretraining on smaller, task-specific datasets.
\end{itemize}

This suggests that \textbf{video prediction is a foundational task} for video understanding, much as next-token prediction is foundational for language. Scaling video prediction to billions of hours of video and trillion-parameter models remains an active frontier, with potential to unlock powerful open-domain video AI.

\begin{notebox}
\textbf{Key Takeaway}: Next-frame prediction mirrors language modelling's philosophy—leverage the structure of raw unlabelled data to learn powerful, transferable representations. For video, this means learning dynamics, causality, and scene understanding as an emergent property of predicting the sensory signal.
\end{notebox}
